% !TeX root = ../main.tex
\section{P2. Konkurensi, Availability, dan Degradasi}
\label{sec:p2}
\subsection{Model konkurensi dan isolasi}
Go menjalankan worker BMKG dan PVMBG secara independen. Dua endpoint BMKG di-fetch paralel, lalu hasilnya diterapkan secara berurutan dalam worker BMKG. Worker tidak memulai siklus berikutnya sebelum siklus aktif selesai; konkurensi fetch tidak tumbuh tanpa batas. Tiap endpoint mempunyai state circuit breaker sendiri.

Jalur baca client mengambil data dari Canonical Store melalui query internal Aggregator, sehingga tidak ikut menunggu PVMBG. Pool PostgreSQL query dipisahkan dari pool ingest. Pemisahan pool membatasi kompetisi koneksi, tetapi keduanya masih memakai proses dan database yang sama. Pemakaian CPU, disk, dan lock database karena itu belum sepenuhnya terisolasi.
\diagram{03-query-isolation}{Pemisahan request baca dan polling sumber. Pembacaan seismik tetap menggunakan store ketika panggilan PVMBG lambat atau gagal.}{fig:query}

Cabang paralel pada Gambar~\ref{fig:query} menunjukkan dua pekerjaan yang tidak saling menunggu. Timeout dan circuit breaker berlaku pada worker polling, sedangkan permintaan client mempunyai batas waktu query sendiri. Pemeriksaan scope dan batas beban dilakukan sebelum query agar permintaan yang tidak berhak atau melampaui kapasitas tidak memakai koneksi database. Proyeksi diterapkan setelah query untuk membentuk respons sesuai hak client. Ketika PVMBG gagal, jalur baca tetap dapat mengembalikan data tersimpan beserta status basi; independensi ini tidak berlaku jika Canonical Store sendiri tidak tersedia.

\begin{tabularx}{\linewidth}{P{0.25\linewidth}Y}
\toprule Kegagalan & Mekanisme dan batasnya\\\midrule
Head-of-line blocking & Worker sumber terpisah dan materialisasi data mengeluarkan HTTP sumber dari jalur request client. Read tidak menjamin data terbaru sensor.\\
Resource exhaustion & Semaphore request aktif, token bucket per identitas, pagination, pool koneksi terbatas, dan timeout. Penolakan terkendali dilaporkan sebagai 429, bukan disamarkan sebagai sukses.\\
Tidak ada degradasi & Status per endpoint diagregasi menjadi status sumber. Data lama tetap tersedia dengan \texttt{stale\_since}; sumber tanpa data usable dapat menghasilkan 503 eksplisit.\\
\bottomrule
\end{tabularx}

\subsection{Konfigurasi dasar}
\begin{tabularx}{\linewidth}{P{0.50\linewidth}Y}
\toprule Parameter & Nilai default yang dilaporkan\\\midrule
Interval polling BMKG dan PVMBG & BMKG 2 s dan PVMBG 5 s; jitter kurang dari 10\% interval.\\
Overlap watermark & 10 s per endpoint.\\
Timeout HTTP sumber BMKG dan PVMBG & BMKG 1 s dan PVMBG 4 s.\\
Timeout operasi dan transaksi PostgreSQL & 2 s; deadline parent yang lebih pendek tetap berlaku.\\
Timeout client-api ke Aggregator & Maksimum 1,5 s total.\\
Breaker & Buka setelah 3 kegagalan fetch; cooldown 10 s.\\
Pool PostgreSQL ingest dan query & Ingest 5 koneksi dan query 3 koneksi.\\
Request aktif client-api & Maksimum 100.\\
Token bucket per identitas & 100 request per detik; kapasitas burst 200.\\
Ukuran halaman & Default 100, maksimum 500.\\
\bottomrule
\end{tabularx}
\coderef{services/aggregator/internal/config/config.go}
\coderef{services/client-api/internal/config/config.go}

Breaker menahan fetch baru selama sumber berulang kali gagal. Poll yang dilewati bukan percobaan HTTP baru. Saat cooldown selesai, percobaan yang berhasil memulihkan alur tanpa restart. Status sumber sehat baru dapat diketahui setelah polling dan pembaruan status berikutnya; mematikan outage tidak membuat metadata langsung HEALTHY pada detik yang sama.

\subsection{Kebaruan data dan kontrak galat}
\texttt{stale\_since} menunjukkan awal masalah yang terdeteksi, bukan waktu persis instansi mulai mati. Pada kondisi normal, perkiraan lag mencakup waktu menunggu polling, jitter, durasi fetch, dan transaksi penyimpanan. Secara konseptual,
\[
L_{\mathrm{ingest}} \approx W_{\mathrm{poll}}+J+T_{\mathrm{HTTP}}+T_{\mathrm{tx}}.
\]
Perkiraan tersebut tidak menjadi jaminan batas waktu layanan. Siklus yang panjang, kegagalan database, ketidakselarasan jam, atau circuit breaker dapat menambah waktu ingest. Outage berkepanjangan membuat umur data terus bertambah. Untuk jalur consumer, tambahkan waktu tunggu relay, ACK broker, dan pemrosesan consumer. Klaim ``maksimal satu interval polling'' tidak digunakan.

Filter tidak valid menghasilkan 400, identitas atau token invalid 401, akses field terlarang 403, detail yang tidak ditemukan 404, beban ditolak 429, dan dependensi tidak tersedia 503. Hasil filter kosong yang sah tetap 200 dengan array kosong. Error tidak membocorkan detail SQL atau secret.

Metadata \texttt{sources} tersedia pada list serta detail biasa dan raw. Metadata tetap disertakan ketika client memilih sebagian field hazard, sehingga pembaca detail dapat mengetahui status basi. Instrumentasi outbound mencatat latency dan correlation ID untuk HTTP sumber, PostgreSQL, Redis, serta Kafka tanpa menulis SQL, argumen perintah, atau secret. Permintaan protokol Kafka yang mencakup beberapa record memakai identitas operasi tersendiri; publish dan commit per record mempertahankan correlation ID event.

\subsection{Konfigurasi dan hasil pengukuran}
Pengukuran terbaru dilakukan pada 8 Oktober 2026 terhadap revisi \texttt{2353479} menggunakan k6 0.57.0 native Windows. Docker Desktop 28.0.4 menyediakan 16 CPU dan 7.936.241.664 byte RAM untuk lingkungan container. Tool load dibangun dengan Go 1.23.6; service memakai Go 1.24.2. Angka di bagian ini diambil dari artefak JSON pengujian ulang. Pengukuran 6 Oktober tetap tersedia sebagai bukti historis dan tidak ditimpa.

PVMBG diatur dengan delay minimum dan maksimum 3 s. Skenario paralel menjalankan 10 VU seismic dan 10 VU volcanic selama 60 s. Skenario sustained menjalankan 50 VU selama 90 s dengan jeda 0,1 s per iterasi; setiap VU memiliki sesi refresh sendiri. Ini model closed loop. Semua VU menggunakan identitas Tim Lapangan yang sama, sehingga berbagi rate limit client. Skenario outage memakai 5 VU per fase (20 s outage, 30 s recovery) dan satu VU kontrol.

\begin{table}[H]\centering\small
\begin{tabular}{lrrrrr}\toprule
Skenario & Bisnis/s & Sukses/s & p50 (ms) & p95 (ms) & p99 (ms)\\\midrule
seismic-only & 188,49 & 99,72 & 6,63 & 12,27 & 18,67\\
sustained & 471,80 & 97,59 & 6,89 & 13,93 & 27,11\\
outage & 32,47 & 32,47 & 4,56 & 6,82 & 7,91\\
\bottomrule\end{tabular}
\caption{Pengujian ulang k6 pada 8 Oktober 2026. Latency hanya HTTP 200; bisnis/s mencakup 429.}\end{table}


Percentile dalam tabel adalah respons HTTP 200 gabungan masing-masing skenario. Untuk kriteria khusus BMKG-only, p95 \textbf{seismic saja 10,92 ms}, lebih kecil dari 300 ms ketika pembacaan volcanic berlangsung bersamaan. Durasi mengikuti definisi k6 yang mencakup waktu pengiriman, penantian respons, dan penerimaan data, tetapi tidak mencakup pembentukan koneksi~\cite{k6-metrics}. Penolakan 429 tidak dimasukkan ke distribusi latency sukses.

Pada sustained, 42.801 respons bisnis terdiri atas 9.028 HTTP 200 dan 33.773 HTTP 429. Ada 10 penolakan autentikasi 429 yang dicatat terpisah; auth error setelah mekanisme retry 0\%. Throughput sebesar 475,03 respons bisnis per detik mencakup penolakan akibat pembatasan beban. Throughput HTTP 200 sebesar 100,20 respons per detik konsisten dengan batas 100 request per detik untuk satu identitas yang dipakai seluruh VU. Sebanyak 78,91\% request bisnis ditolak secara terkontrol. Hasil ini menunjukkan kestabilan pada pola tersebut, tetapi belum mengukur kapasitas maksimum server.

Error bisnis dihitung sebagai jumlah respons gagal di luar penolakan terkontrol dibagi jumlah respons bisnis selain 429. Pada sustained hasilnya $0/9028=0\%$. Sampel 429 dikeluarkan dari pembilang \emph{dan} penyebut agar tidak mengencerkan error rate. Total HTTP bawaan k6 mencampurkan auth, control, dan retry, sehingga bukan metrik throughput bisnis utama.

\diagram{06-load-results}{Grafik yang dibangkitkan langsung dari JSON hasil k6 dan sampling koneksi. Percentile memakai respons sukses; koneksi diukur terpisah dari jumlah VU.}{fig:load}

Runner membaca \nolinkurl{/proc/net/tcp} dan \nolinkurl{/proc/net/tcp6} sekitar setiap satu detik untuk menghitung koneksi TCP ESTABLISHED pada port client-api. Sebanyak 71 sampel menunjukkan sedikitnya 50 koneksi teramati terus-menerus selama \textbf{88,53 s}. Koneksi aktif secara TCP tidak berarti semua request sedang diproses serentak; jumlah VU saja juga tidak dipakai sebagai bukti koneksi.

Selama outage, seluruh 480 respons seismic tetap sukses dengan data dan status BMKG HEALTHY. Seluruh 480 respons volcanic mengembalikan data tersimpan dengan penanda basi. Pada fase recovery, 420 dari 730 respons (57,53\%) sudah kembali HEALTHY; 310 respons lainnya masih berada dalam masa pemulihan, tetapi tetap HTTP 200. Polling sumber kembali sukses pada 15:46:56 WIB sebelum runner mengembalikan konfigurasi mock pada 15:47:15 WIB. Pemulihan tidak memerlukan restart aplikasi. Persentase recovery ini memenuhi threshold internal script di atas 50\%, tetapi bukan jaminan waktu pemulihan seketika.

Metrik dan konfigurasi lengkap tersimpan pada \reliabilityevidence{README.md}. Angka berasal dari \reliabilityevidence{load/seismic-only.json}, \reliabilityevidence{load/sustained.json}, \reliabilityevidence{load/connections.json}, serta \reliabilityevidence{load/outage.json}.

\subsection{Validitas hasil dan alternatif}
Percobaan awal k6 di container menghasilkan durasi negatif, misalnya minimum -541,10 ms. Walaupun threshold k6 saat itu tidak gagal, hasil tersebut tidak digunakan untuk klaim performa. Artefaknya tetap disimpan pada \evidence{load-docker-invalid/seismic-only.json}. Penyebab lingkungan belum diisolasi secara pasti. Runner kini menolak nilai durasi negatif, dan seluruh pengukuran final diulang dengan k6 native Windows terhadap stack yang sama.

Alternatif request-time fan-out menjaga kemungkinan kebaruan lebih tinggi, tetapi menempatkan latensi dan kegagalan sumber di jalur client. Materialisasi dipilih karena data historis masih berguna dan P2 menuntut degradasi terkendali. Cache dan micro-cache dapat mengurangi beban query, tetapi menambah invalidasi dan state; belum digunakan. Banyak instance dengan load balancer dapat meningkatkan kapasitas, tetapi memerlukan limit terdistribusi dan koordinasi penulis yang berada di luar M1. Hasil satu mesin ini tidak digeneralisasi sebagai kapasitas produksi.
